% Options for packages loaded elsewhere
\PassOptionsToPackage{unicode}{hyperref}
\PassOptionsToPackage{hyphens}{url}
\documentclass[
]{article}
\usepackage{xcolor}
\usepackage{amsmath,amssymb}
\setcounter{secnumdepth}{-\maxdimen} % remove section numbering
\usepackage{iftex}
\ifPDFTeX
  \usepackage[T1]{fontenc}
  \usepackage[utf8]{inputenc}
  \usepackage{textcomp} % provide euro and other symbols
\else % if luatex or xetex
  \usepackage{unicode-math} % this also loads fontspec
  \defaultfontfeatures{Scale=MatchLowercase}
  \defaultfontfeatures[\rmfamily]{Ligatures=TeX,Scale=1}
\fi
\usepackage{lmodern}
\ifPDFTeX\else
  % xetex/luatex font selection
\fi
% Use upquote if available, for straight quotes in verbatim environments
\IfFileExists{upquote.sty}{\usepackage{upquote}}{}
\IfFileExists{microtype.sty}{% use microtype if available
  \usepackage[]{microtype}
  \UseMicrotypeSet[protrusion]{basicmath} % disable protrusion for tt fonts
}{}
\makeatletter
\@ifundefined{KOMAClassName}{% if non-KOMA class
  \IfFileExists{parskip.sty}{%
    \usepackage{parskip}
  }{% else
    \setlength{\parindent}{0pt}
    \setlength{\parskip}{6pt plus 2pt minus 1pt}}
}{% if KOMA class
  \KOMAoptions{parskip=half}}
\makeatother
\usepackage{longtable,booktabs,array}
\usepackage{calc} % for calculating minipage widths
% Correct order of tables after \paragraph or \subparagraph
\usepackage{etoolbox}
\makeatletter
\patchcmd\longtable{\par}{\if@noskipsec\mbox{}\fi\par}{}{}
\makeatother
% Allow footnotes in longtable head/foot
\IfFileExists{footnotehyper.sty}{\usepackage{footnotehyper}}{\usepackage{footnote}}
\makesavenoteenv{longtable}
\usepackage{graphicx}
\makeatletter
\newsavebox\pandoc@box
\newcommand*\pandocbounded[1]{% scales image to fit in text height/width
  \sbox\pandoc@box{#1}%
  \Gscale@div\@tempa{\textheight}{\dimexpr\ht\pandoc@box+\dp\pandoc@box\relax}%
  \Gscale@div\@tempb{\linewidth}{\wd\pandoc@box}%
  \ifdim\@tempb\p@<\@tempa\p@\let\@tempa\@tempb\fi% select the smaller of both
  \ifdim\@tempa\p@<\p@\scalebox{\@tempa}{\usebox\pandoc@box}%
  \else\usebox{\pandoc@box}%
  \fi%
}
% Set default figure placement to htbp
\def\fps@figure{htbp}
\makeatother
\setlength{\emergencystretch}{3em} % prevent overfull lines
\providecommand{\tightlist}{%
  \setlength{\itemsep}{0pt}\setlength{\parskip}{0pt}}
\usepackage{bookmark}
\IfFileExists{xurl.sty}{\usepackage{xurl}}{} % add URL line breaks if available
\urlstyle{same}
\hypersetup{
  pdftitle={Causal Nonseparability and Environmental Objectivity},
  hidelinks,
  pdfcreator={LaTeX via pandoc}}

\title{\textbf{Causal Nonseparability and Environmental Objectivity}}
\author{}
\date{}

\begin{document}
\maketitle

\emph{Independence, Control-Assisted Records, and Quantum-Switch
Benchmarks}

\textbf{Scott Muniz}

Dr. Ware Technology Services \textbar{} scott@drware.com

\emph{Research Article \textbar{} Version 5 \textbar{} Final
manuscript-level validation update \textbar{} August 2026}

\section{Popular Summary}\label{popular-summary}

The environment can act like a witness: many separate pieces of it may
carry the same information about a quantum system. Quantum Darwinism
explains how such repeated records support an objective classical
description. A different line of research asks what happens when a
quantum system controls whether operation A occurs before B or B before
A. We connect these subjects without assuming that either one causes the
other. Objective records and indefinite causal order are independent
properties. The order-control system can contain extra information, but
only when it is measured together with an environmental fragment, and
the extra information obeys a strict budget. Ignoring the control
removes the coherence-dependent contribution. The results separate
redundant environmental records, joint measurements of the environment
and order control, and independent certification of causal structure,
while providing a four-qubit protocol with testable predictions.

\section{Abstract}\label{abstract}

Environmental objectivity and indefinite causal order concern distinct
operational structures: redundant pointer information in environmental
fragments and the convex causal structure of a process. We prove their
logical independence in a composite class containing fixed-order record
circuits and the quantum switch. We then establish exact equal-prior
Helstrom neutrality for a binary record model with arbitrary mismatch
and unequal interaction strengths: coherent order adds two null
directions to the conditional-difference spectrum without changing its
nonzero singular values. This exact no-advantage result coexists with a
resource-matched joint-access Holevo advantage of 0.083977 bit. For a
classical pointer label Λ, order control T, and fragment F, we derive a
nonnegative access decomposition separating fragment information,
classical branch-label information, and coherence-sensitive information.
In irreducibly coupled controlled-order dynamics, exact unconditional
pointer preservation is equivalent to preservation in each orthogonal
branch. An exact postselection identity and probability-weighted bounds
are controlled by I(Λ:T\textbar F). Many-fragment calculations, local
spectrum-broadcast diagnostics, collision-noise surfaces, exact
gate-level statevector simulation, reduced-normal-form checks, and
finite-shot ideal sampling delimit the capability. Finally,
deterministic local record augmentation of a causally nonseparable
process preserves causal nonseparability, and every witness for the base
process pulls back through the discard map to a witness for the
extension. The resource is therefore bounded, task dependent, and access
dependent rather than a universal enhancement of objectivity.

\section{I. Introduction}\label{i.-introduction}

Quantum Darwinism describes operational objectivity through redundant
environmental records of a system pointer {[}1-8{]}. Process-matrix and
higher-order-operation frameworks instead classify transformations and
processes that cannot be reduced to a definite or probabilistic causal
order {[}9-13,19{]}. These frameworks answer different questions.
Objectivity asks which information is available in multiple disjoint
fragments. Causal nonseparability asks whether the higher-order process
connecting local laboratories lies outside the causally separable cone.

This work establishes a disciplined relation between the two frameworks.
Neither structure implies the other. When the same interactions generate
records and participate in coherently controlled orders, branchwise
pointer compatibility controls deterministic preservation, whereas
access to the control supplies a bounded information resource. The
resource is task dependent: it can increase joint Holevo information,
leave equal-prior minimum-error discrimination exactly unchanged, and
fail to improve environmental redundancy.

Why this matters now. Quantum-switch experiments and certification
methods have advanced rapidly {[}15,21-23{]}, while recent theoretical
and experimental work has sharpened the conditions for Quantum
Darwinism, strong objectivity, and spectrum-broadcast structure
{[}17,18,20,24{]}. A unified operational treatment is needed to prevent
information gain, redundant record formation, and causal certification
from being treated as interchangeable. The present framework supplies
that separation and a small-register route for testing it.

The principal PRX Quantum case is an exceptional connection between
environmental objectivity and indefinite causal order. The connection is
not a claim that indefinite order universally improves objectivity.
Rather, it identifies exactly which conclusions survive when access is
restricted, when outcomes are postselected, and when causal structure is
certified independently.

The main contributions are: (i) a logical-independence theorem; (ii) an
exact equal-prior Helstrom-neutrality theorem; (iii) an exact
quantum-control access decomposition; (iv) branchwise preservation
theorems; (v) a universal postselection information budget; (vi)
many-fragment, noise, and local spectrum-broadcast diagnostics; (vii) a
four-qubit validation protocol; and (viii) an inherited causal-witness
theorem for deterministic local record extensions.

Scope is deliberately bounded. We do not claim an advantage over the
full causally separable process cone, a universal enhancement of
objectivity, a global distance-to-spectrum-broadcast theorem, or
physical-device validation.

\textbf{Table I. Theorem roadmap and logical architecture.}

\begin{longtable}[]{@{}
  >{\raggedright\arraybackslash}p{(\linewidth - 6\tabcolsep) * \real{0.2500}}
  >{\raggedright\arraybackslash}p{(\linewidth - 6\tabcolsep) * \real{0.2500}}
  >{\raggedright\arraybackslash}p{(\linewidth - 6\tabcolsep) * \real{0.2500}}
  >{\raggedright\arraybackslash}p{(\linewidth - 6\tabcolsep) * \real{0.2500}}@{}}
\toprule\noalign{}
\begin{minipage}[b]{\linewidth}\raggedright
\textbf{Layer}
\end{minipage} & \begin{minipage}[b]{\linewidth}\raggedright
\textbf{Question}
\end{minipage} & \begin{minipage}[b]{\linewidth}\raggedright
\textbf{Core result}
\end{minipage} & \begin{minipage}[b]{\linewidth}\raggedright
\textbf{Location}
\end{minipage} \\
\midrule\noalign{}
\endhead
\bottomrule\noalign{}
\endlastfoot
Independence & Do objectivity and causal nonseparability imply one
another? & No; all four logical sectors are constructively populated. &
Theorem 1 \\
Decision task & Does coherent order improve equal-prior binary
discrimination? & No in the declared binary model; the trace norm is
exactly unchanged. & Sec. VI \\
Access and preservation & Where can extra information reside, and when
is the pointer preserved? & Nonnegative access decomposition and
branchwise preservation criteria. & Theorems 2-4 \\
Operational limits & Can measurement or postselection create
unrestricted record gain? & No; average and selected gains obey
probability-weighted budgets. & Theorems 5-6 \\
Certification and test & Does record augmentation retain causal
certification? & Yes under deterministic local recovery; the witness
pulls back. & Theorem 7 and Sec. IX \\
\end{longtable}

\section{II. Operational setting and
notation}\label{ii.-operational-setting-and-notation}

\subsection{A. Global assumptions}\label{a.-global-assumptions}

All Hilbert spaces are finite dimensional. Logarithms are base two, so
information quantities are measured in bits. The trace norm is
\textbar\textbar X\textbar\textbar₁=Tr√(X†X). Unless otherwise stated,
the binary decision problem uses equal priors. Classical registers are
represented by orthogonal basis states, and conditional mutual
information is I(A:B\textbar C)=S(AC)+S(BC)-S(C)-S(ABC).

\textbf{Table II. Principal notation and conventions.}

\begin{longtable}[]{@{}
  >{\raggedright\arraybackslash}p{(\linewidth - 2\tabcolsep) * \real{0.5000}}
  >{\raggedright\arraybackslash}p{(\linewidth - 2\tabcolsep) * \real{0.5000}}@{}}
\toprule\noalign{}
\begin{minipage}[b]{\linewidth}\raggedright
\textbf{Symbol}
\end{minipage} & \begin{minipage}[b]{\linewidth}\raggedright
\textbf{Meaning}
\end{minipage} \\
\midrule\noalign{}
\endhead
\bottomrule\noalign{}
\endlastfoot
Λ & Classical pointer label with prior qλ \\
T & Quantum order-control system \\
F & Accessible environmental fragment \\
E₁,E₂ & Record subsystems written by the two local interactions \\
U\_A,U\_B & Record-generating local unitaries \\
Qλ & Pointer-sector projector \\
D\_T & Complete dephasing of T in the order-branch basis \\
Y & Classical outcome register of an instrument on T \\
W\_sw,W\_ext & Base switch process and record-augmented process \\
R,R* & Deterministic local recovery/discard map and its Hilbert-Schmidt
adjoint \\
S\_sw,S\_ext & Base causal witness and its pulled-back extension \\
\end{longtable}

\subsection{B. Pointer labels, fragments, and
objectivity}\label{b.-pointer-labels-fragments-and-objectivity}

Let Λ be a classical pointer label with prior qλ and conditional states
ρ\_\{TF\textbar λ\}. The associated classical-quantum state is

ρ\_\{ΛTF\} = Σλ qλ \textbar λ\textgreater\textless λ\textbar{} ⊗
ρ\_\{TF\textbar λ\}. (1)

For a fragment F, I(Λ:F) equals the Holevo information of the
conditional ensemble. Operational redundancy at deficit δ requires many
disjoint fragments to carry at least (1-δ)H(Λ). Ordinary redundancy is
not identified with strong Quantum Darwinism or spectrum-broadcast
structure without their additional conditions; in particular, strong
Quantum Darwinism and strong independence provide conditions connected
to spectrum-broadcast structure {[}20{]}.

\subsection{C. Causal separability and composite process
class}\label{c.-causal-separability-and-composite-process-class}

For the tripartite switch setting with a final party and no final
output, a causally separable process admits a convex decomposition

W\_sep = q W\^{}(A≺B≺C) + (1-q) W\^{}(B≺A≺C). (2)

A causal witness is nonnegative on every causally separable process and
negative on at least one causally nonseparable process {[}10-12{]}.
Record quality and causal certification are evaluated separately.

\begin{longtable}[]{@{}
  >{\raggedright\arraybackslash}p{(\linewidth - 0\tabcolsep) * \real{1.0000}}@{}}
\toprule\noalign{}
\begin{minipage}[b]{\linewidth}\raggedright
\textbf{Definition 1. Composite record-process class.\\
}Let C\_rec contain (a) fixed-order processes with deterministic local
record channels, (b) the unaltered quantum switch without appended
records, and (c) record-augmented switch processes. The class is closed
under adjoining deterministic local record channels and under
deterministic local discard of those records. The operations preserve
the validity constraints of the underlying process space.\strut
\end{minipage} \\
\midrule\noalign{}
\endhead
\bottomrule\noalign{}
\endlastfoot
\end{longtable}

\section{III. Independence and control-access
decomposition}\label{iii.-independence-and-control-access-decomposition}

\subsection{Theorem 1. Logical
independence}\label{theorem-1.-logical-independence}

Within C\_rec, environmental objectivity is neither necessary nor
sufficient for causal nonseparability.

Proof. A fixed-order process with compatible repeated records is
objective but causally separable. The unaltered causally nonseparable
quantum switch, or a switch with blank or incompatible record channels,
is nonobjective but causally nonseparable. A switch with
common-pointer-preserving record channels can satisfy both properties. A
fixed-order nonrecording process satisfies neither. The quantum switch
is treated as a higher-order controlled-order process as introduced in
Ref. {[}19{]}, and causal nonseparability is certified independently
through the process framework {[}10-12,15{]}. These constructions
populate all four sectors. □

\textbf{Table III. Constructive population of the four logical sectors.}

\begin{longtable}[]{@{}
  >{\raggedright\arraybackslash}p{(\linewidth - 4\tabcolsep) * \real{0.3333}}
  >{\raggedright\arraybackslash}p{(\linewidth - 4\tabcolsep) * \real{0.3333}}
  >{\raggedright\arraybackslash}p{(\linewidth - 4\tabcolsep) * \real{0.3333}}@{}}
\toprule\noalign{}
\begin{minipage}[b]{\linewidth}\raggedright
\textbf{Environmental objectivity}
\end{minipage} & \begin{minipage}[b]{\linewidth}\raggedright
\textbf{Causal nonseparability}
\end{minipage} & \begin{minipage}[b]{\linewidth}\raggedright
\textbf{Construction}
\end{minipage} \\
\midrule\noalign{}
\endhead
\bottomrule\noalign{}
\endlastfoot
Yes & No & Fixed order with compatible repeated records \\
No & Yes & Unaltered switch or switch with blank/incompatible records \\
Yes & Yes & Switch with common-pointer-preserving record channels \\
No & No & Fixed-order nonrecording process \\
\end{longtable}

\subsection{Theorem 2. Nonnegative access
decomposition}\label{theorem-2.-nonnegative-access-decomposition}

Let D\_T completely dephase T in the order-branch basis. Define
I\_F=I(Λ:F)\_ρ, I\_class=I(Λ:TF)\_\{D\_T(ρ)\}, and I\_coh=I(Λ:TF)\_ρ.
Then

I\_coh - I\_F = Δ\_coh + Δ\_label, Δ\_coh ≥ 0, Δ\_label ≥ 0, (3)

where Δ\_coh=I\_coh-I\_class and Δ\_label=I\_class-I\_F.

Proof. D\_T is a channel local to T. Data processing therefore gives
I(Λ:TF)\_ρ ≥ I(Λ:TF)\_\{D\_T(ρ)\}, proving Δ\_coh≥0. Because D\_T acts
only on T and F is disjoint, dephasing commutes with tracing out T;
hence the ΛF marginal is unchanged. The chain rule applied to the
dephased state gives

I(Λ:TF)\_\{D\_T(ρ)\} = I(Λ:F)\_\{D\_T(ρ)\} +
I(Λ:T\textbar F)\_\{D\_T(ρ)\}.

The unchanged ΛF marginal implies I(Λ:F)\_\{D\_T(ρ)\}=I(Λ:F)\_ρ.
Therefore

Δ\_label = I(Λ:T\textbar F)\_\{D\_T(ρ)\} ≥ 0. (4)

The split is defined relative to the declared branch-basis dephasing and
access structure; it is not asserted to be a unique partial-information
decomposition. □

\section{IV. Controlled-order record
dynamics}\label{iv.-controlled-order-record-dynamics}

Let U\_A record the Z pointer into E₁ and U\_B record a θ-rotated
pointer into E₂. The controlled-order unitary is

U\_ord = \textbar0\textgreater\textless0\textbar\_T ⊗ U\_B U\_A +
\textbar1\textgreater\textless1\textbar\_T ⊗ U\_A U\_B. (5)

\begin{longtable}[]{@{}
  >{\raggedright\arraybackslash}p{(\linewidth - 0\tabcolsep) * \real{1.0000}}@{}}
\toprule\noalign{}
\begin{minipage}[b]{\linewidth}\raggedright
\textbf{Definition 2. Exact pointer preservation.\\
}An operator V preserves the pointer projectors \{Qλ\} exactly when
{[}V,Qλ{]}=0 for every λ. This operator-level condition is equivalent to
block-diagonality in the pointer decomposition and prevents transitions
between pointer sectors. For a state supported in one sector, it implies
nondemolition preservation of the corresponding pointer value.\strut
\end{minipage} \\
\midrule\noalign{}
\endhead
\bottomrule\noalign{}
\endlastfoot
\end{longtable}

\subsection{Theorem 3. Common-pointer
sufficiency}\label{theorem-3.-common-pointer-sufficiency}

If {[}U\_A,Qλ{]}={[}U\_B,Qλ{]}=0 for every pointer projector Qλ, then
both definite orders, U\_ord, and every nonzero postselected operator
Kφ=\textless φ\textbar U\_ord\textbar τ\textgreater{} preserve Qλ.
Products and linear combinations of operators in the commutant remain in
the commutant.

\subsection{Theorem 4. Orthogonal-branch
necessity}\label{theorem-4.-orthogonal-branch-necessity}

For an initial control state τ\_T with branch weights
p₀=\textless0\textbar τ\_T\textbar0\textgreater{} and
p₁=\textless1\textbar τ\_T\textbar1\textgreater, exact unconditional
operator-level preservation by U\_ord holds if and only if both order
products preserve the pointer projector, because

{[}U\_ord,I\_T⊗Qλ{]} = \textbar0\textgreater\textless0\textbar⊗{[}U\_B
U\_A,Qλ{]} + \textbar1\textgreater\textless1\textbar⊗{[}U\_A U\_B,Qλ{]}.
(6)

The terms occupy orthogonal control blocks, so cross-order coherence
cannot cancel operator-level leakage. For a projector Qλ, define the
nonnegative Hilbert-Schmidt leakage
Lλ(V)=\textbar\textbar(I-Qλ)VQλ\textbar\textbar₂². With initial control
branch weights p₀ and p₁, tracing out T gives Lλ(U\_ord)=p₀Lλ(U\_B
U\_A)+p₁Lλ(U\_A U\_B). Thus zero unconditional leakage with both weights
nonzero is equivalent to zero leakage in each occupied branch. If one
branch has zero weight, necessity applies only to the occupied branch.
The same conclusion holds for any nonnegative leakage functional
additive on orthogonal control blocks.

\section{V. Postselection information
budget}\label{v.-postselection-information-budget}

Let an instrument \{M\_y\} on T store its accessible outcome in a
classical register Y. The post-instrument state is ρ\_\{ΛFY\}=Σ\_y p\_y
ρ\_\{ΛF\textbar y\}⊗\textbar y\textgreater\textless y\textbar, where
p\_y\textgreater0 for every selected outcome retained in the formulas.
Define χ\_y=I(Λ:F)\_\{ρ\textbar y\} and χ=I(Λ:F)\_ρ.

\subsection{Theorem 5. Exact average
identity}\label{theorem-5.-exact-average-identity}

Σ\_y p\_y χ\_y - χ = I(Λ:Y\textbar F) - I(Λ:Y). (7)

Proof. Because Y is classical, I(Λ:F\textbar Y)=Σ\_y p\_y
I(Λ:F)\_\{ρ\textbar y\}=Σ\_y p\_yχ\_y. The two chain-rule expansions

I(Λ:FY)=I(Λ:Y)+I(Λ:F\textbar Y)=I(Λ:F)+I(Λ:Y\textbar F)

imply Eq. (7) after rearrangement. □

\subsection{Theorem 6. Average and selected-outcome
bounds}\label{theorem-6.-average-and-selected-outcome-bounds}

Σ\_y p\_y(χ\_y-χ) ≤ I(Λ:T\textbar F), (8)

p\_y(χ\_y-χ) ≤ I(Λ:T\textbar F) + (1-p\_y)χ. (9)

Proof. The measurement channel T→Y and conditional data processing give
I(Λ:Y\textbar F)≤I(Λ:T\textbar F). Since I(Λ:Y)≥0, Eq. (7) yields Eq.
(8). For a chosen y, write the average deviation as
p\_y(χ\_y-χ)+Σ\_\{z≠y\}p\_z(χ\_z-χ). Because χ\_z≥0, every complementary
deviation is at least -χ, and hence Σ\_\{z≠y\}p\_z(χ\_z-χ)≥-(1-p\_y)χ.
Combining this lower bound with Eq. (8) gives Eq. (9). Thus
postselection redistributes fragment information; a favorable outcome
must be reported with its probability and complementary outcomes. □

\section{VI. Exact equal-prior decision
neutrality}\label{vi.-exact-equal-prior-decision-neutrality}

This section treats the exact binary decision task. For equal priors,
let Δ\_AB=ρ\^{}\{AB\}\_\{E₂\textbar0\}-ρ\^{}\{AB\}\_\{E₂\textbar1\},
Δ\_BA=ρ\^{}\{BA\}\_\{E₂\textbar0\}-ρ\^{}\{BA\}\_\{E₂\textbar1\}, and
Δ\_coh=ρ\^{}\{coh\}\_\{TE₂\textbar0\}-ρ\^{}\{coh\}\_\{TE₂\textbar1\}.
The fixed-order operators are 2×2 Hermitian matrices on E₂; the coherent
operator is a 4×4 Hermitian matrix on T⊗E₂. We evaluate them in the
ordered bases \{\textbar0\textgreater,\textbar1\textgreater\} and
\{\textbar0\_T0\_E₂\textgreater,\textbar0\_T1\_E₂\textgreater,\textbar1\_T0\_E₂\textgreater,\textbar1\_T1\_E₂\textgreater\}.

For the declared binary record model, the symbolic certificate
constructs the physical-basis operators and verifies a unitary reduction
to the following spectral normal forms. Let s=sin(α\_B/2)cosθ. There
exist unitaries V\_AB, V\_BA, and V\_coh such that

V\_AB Δ\_AB V\_AB† = diag(s,-s), V\_BA Δ\_BA V\_BA† = diag(s,-s),

V\_coh Δ\_coh V\_coh† = diag(s,-s,0,0).

These displayed matrices are the exact reduced normal forms, not assumed
representatives. The archived symbolic calculation verifies coefficient
by coefficient that p\_coh(λ)-λ²p\_AB(λ)=0 and p\_coh(λ)-λ²p\_BA(λ)=0
under the trigonometric normalization identities. Consequently,

p\_fix(λ)=det(λI-Δ\_AB)=det(λI-Δ\_BA)=λ²-sin²(α\_B/2)cos²θ. (10)

The coherent matrix has rank two. Elementary block reduction or a
symbolic determinant expansion yields

p\_coh(λ)=det(λI-Δ\_coh)=λ²{[}λ²-sin²(α\_B/2)cos²θ{]}. (11)

Thus Δ\_coh adds two null directions and has the same nonzero
eigenvalues ±\textbar sin(α\_B/2)cosθ\textbar{} as either fixed order.
The equality holds for arbitrary mismatch θ and arbitrary interaction
strengths α\_A and α\_B within the model. The absence of α\_A reflects
the fact that, for this marginal binary decision observable, the first
record interaction changes the embedding but not the two nonzero
singular values. Therefore

\textbar\textbar Δ\_coh\textbar\textbar₁=\textbar\textbar Δ\_AB\textbar\textbar₁=\textbar\textbar Δ\_BA\textbar\textbar₁=2\textbar sin(α\_B/2)cosθ\textbar,
(12)

P\_H=1/2+1/4\textbar\textbar Δ\textbar\textbar₁=1/2+1/2\textbar sin(α\_B/2)cosθ\textbar.
(13)

This equality is not a generic statement that coherent order is useless.
Holevo information depends on the spectra of the individual conditional
states and their average, not only on the binary trace norm. Biased
priors replace Δ by π₀ρ₀-π₁ρ₁, introducing a noncanceling common-state
term; consequently the factorization in Eq. (11) need not hold. The
result is therefore an exact, task-specific no-advantage theorem.

\section{VII. Numerical and operational
results}\label{vii.-numerical-and-operational-results}

\subsection{A. Matched-family joint-access
capability}\label{a.-matched-family-joint-access-capability}

For fixed dimensions, initial states, one use of each operation, and
comparison with the two fixed orders, the model reaches a joint
control-fragment Holevo advantage of 0.0839771759 bit near θ=0.700108
and full record strength. Discarding T removes the cross-order
contribution and produces no fragment-only advantage over the
corresponding fixed-order family.

\textbf{Table IV. Verified benchmark claims carried into the main text.}

\begin{longtable}[]{@{}
  >{\raggedright\arraybackslash}p{(\linewidth - 6\tabcolsep) * \real{0.2500}}
  >{\raggedright\arraybackslash}p{(\linewidth - 6\tabcolsep) * \real{0.2500}}
  >{\raggedright\arraybackslash}p{(\linewidth - 6\tabcolsep) * \real{0.2500}}
  >{\raggedright\arraybackslash}p{(\linewidth - 6\tabcolsep) * \real{0.2500}}@{}}
\toprule\noalign{}
\begin{minipage}[b]{\linewidth}\raggedright
\textbf{Quantity}
\end{minipage} & \begin{minipage}[b]{\linewidth}\raggedright
\textbf{Coherent / joint-access result}
\end{minipage} & \begin{minipage}[b]{\linewidth}\raggedright
\textbf{Matched fixed-order result}
\end{minipage} & \begin{minipage}[b]{\linewidth}\raggedright
\textbf{Conclusion}
\end{minipage} \\
\midrule\noalign{}
\endhead
\bottomrule\noalign{}
\endlastfoot
Benchmark & θ=0.7001079286; α=π; B-family/full-strength comparison &
Same dimensions, states, operation uses, and record resources &
Resource-matched family \\
B-only fragment, fixed order & 0.477538 bit & AB and BA equal for this
composition & Reference value \\
Joint T+B fragments, 1-3 records & 0.557230 bit & 0.477538 bit &
Advantage 0.079692 bit \\
Joint T+all four B records & 0.561515 bit & 0.477538 bit & Advantage
0.083977 bit \\
Environment-only coherent access & No advantage over optimized
fixed-order mixture & Control-inaccessible reduced state is the equal
order mixture & No fragment-only gain \\
Equal-prior Helstrom & P\_H=1/2+\textbar sin(α\_B/2)cosθ\textbar/2 &
Same expression for AB and BA & Exact neutrality \\
Numerical validation & Gate-level NumPy and Qiskit statevector checks;
reduced-normal-form check; finite-shot ideal sampling & Analytic formula
and matched fixed-order circuits & Computational validation; no
physical-device measurements \\
\end{longtable}

Relation to communication-power results. Zhao et al. show that
coherently controlled order can enhance one-shot classical-message
transmission in their communication framework, while also identifying
asymptotic limits relative to shared entanglement {[}22{]}. The present
task is different: it concerns environmental record access and binary
discrimination. Here coherent order increases joint Holevo information
in a matched family but leaves equal-prior Helstrom performance exactly
unchanged and does not improve fragment-only redundancy in the tested
class. The results are complementary and support the conclusion that
coherent-order advantages depend on the operational task, comparator,
and accessible systems.

\subsection{B. Many-fragment
objectivity}\label{b.-many-fragment-objectivity}

Complete subset enumeration through eight records and mismatch-dependent
surfaces for N=2,4,6 show no coherent fragment-only redundancy advantage
over the better fixed order under either mean-fragment or worst-fragment
0.9-bit thresholds. Repeated perfect records form GHZ-type repetition
encodings: proper nonempty subfragments lose the family coherence and
saturate at the one-record ensemble. The complete-family
control-assisted bonus remains bounded, so the bonus per record vanishes
asymptotically.

\textbf{Table V. Many-fragment conclusions. Numerical arrays remain in
the reproducibility package.}

\begin{longtable}[]{@{}
  >{\raggedright\arraybackslash}p{(\linewidth - 6\tabcolsep) * \real{0.2500}}
  >{\raggedright\arraybackslash}p{(\linewidth - 6\tabcolsep) * \real{0.2500}}
  >{\raggedright\arraybackslash}p{(\linewidth - 6\tabcolsep) * \real{0.2500}}
  >{\raggedright\arraybackslash}p{(\linewidth - 6\tabcolsep) * \real{0.2500}}@{}}
\toprule\noalign{}
\begin{minipage}[b]{\linewidth}\raggedright
Domain
\end{minipage} & \begin{minipage}[b]{\linewidth}\raggedright
Mean-fragment result
\end{minipage} & \begin{minipage}[b]{\linewidth}\raggedright
Worst/global result
\end{minipage} & \begin{minipage}[b]{\linewidth}\raggedright
Interpretation
\end{minipage} \\
\midrule\noalign{}
\endhead
\bottomrule\noalign{}
\endlastfoot
2,4,6,8 & Complete subset enumeration at five mismatch values & No
coherent fragment-only improvement over best fixed order & No coherent
fragment-only improvement over best fixed order \\
8 at θ=0.7001079286 & Best fixed: mδ=3, Rδ≈2.67 & Best fixed: mδ=5,
Rδ=1.6 & Coherent fragment-only never reaches 0.9 bit \\
8, one-record mean & Best fixed 0.738769 bit & Coherent fragment-only
0.498335 bit & Joint T+fragment mean 0.642540 bit \\
8, full environment & Best fixed reaches 1 bit by size 5 & Coherent
fragment-only 0.715579 bit & Joint T+all records 0.943735 bit; global
decoder only \\
\end{longtable}

\subsection{C. Local spectrum-broadcast
diagnostics}\label{c.-local-spectrum-broadcast-diagnostics}

At the benchmark mismatch, coherent environment-only singleton records
remain nonorthogonal and at a finite optimized trace distance from every
ideal orthogonal qubit record pair. Conditional total correlations
diagnose additional nonproduct structure. These finite diagnostics do
not constitute a global minimum distance to the full
spectrum-broadcast-state set.

\subsection{D. Noise and phase
structure}\label{d.-noise-and-phase-structure}

Control dephasing removes the coherence-sensitive contribution. Partial
strengths, reduced-state noise, and collision-level target and record
noise leave positive joint-access regions but do not generate a
fragment-only redundancy advantage in the declared tests. Robustness is
regime specific and is not claimed for arbitrary channels.

\section{VIII. Record-augmented causal
certification}\label{viii.-record-augmented-causal-certification}

Let P denote the valid finite-dimensional process-matrix space for the
declared laboratory interfaces. For the base tripartite switch
interface, the Hilbert spaces are ordered A\_I,A\_O,B\_I,B\_O,C\_I,
Charlie has no output, normalized processes have trace 4, and causally
separable processes decompose as W\_sep=W\^{}(A≺B≺C)+W\^{}(B≺A≺C), with
positive semidefinite summands. The A≺B≺C summand satisfies the
replacement-map constraints \_C W=\_\{B\_O C\}W and \_\{B\_I B\_O
C\}W=\_\{A\_O B\_I B\_O C\}W; the opposite order exchanges A and B. Let
W\_sw∈P be causally nonseparable. Let W\_ext append local record outputs
E\_A,E\_B, and define R=D\_A⊗D\_B⊗id\_C, where D\_A and D\_B are
completely positive deterministic local discard maps satisfying
R(W\_ext)=W\_sw. Because deterministic local processing preserves each
fixed-order cone and the normalization constraints, R maps the extended
causally separable cone into the base causally separable cone.

\subsection{Theorem 7. Ancilla stability and witness
pullback}\label{theorem-7.-ancilla-stability-and-witness-pullback}

W\_ext is causally nonseparable. If S\_sw is a causal witness for W\_sw,
then

S\_ext = R*(S\_sw) (14)

is a causal witness for W\_ext. Here R* is the Hilbert-Schmidt
(Heisenberg) adjoint, defined by

Tr{[}R*(S)W{]}=Tr{[}S R(W){]} (15)

for every compatible witness operator S and process W. Linearity of R
implies linearity of R*. For every causally separable extended process
W\_ext\^{}sep, closure under deterministic local operations gives
R(W\_ext\^{}sep) causally separable; hence

Tr{[}S\_ext W\_ext\^{}sep{]}=Tr{[}S\_sw R(W\_ext\^{}sep){]}≥0. (16)

On the target extension,

Tr(S\_ext W\_ext)=Tr{[}S\_sw R(W\_ext){]}=Tr(S\_sw W\_sw)\textless0.
(17)

Therefore S\_ext is nonnegative on the separable cone and negative on
W\_ext. If W\_ext itself were causally separable, closure under R would
make W\_sw causally separable, a contradiction. Thus the extended
process remains causally nonseparable. □

The causal certificate is independent of the pointer and record
certificates. Common-pointer commutators and fragment-state diagnostics
establish record behavior; the pulled-back witness establishes causal
nonseparability.

\section{IX. Four-qubit validation
route}\label{ix.-four-qubit-validation-route}

The minimal logical register is T,S,E₁,E₂, with global tensor ordering
T⊗S⊗E₁⊗E₂. Exported physical-basis matrices use lexicographic
computational-basis order \textbar t,s,e₁,e₂⟩, with t the most
significant bit. Fixed AB, fixed BA, coherent-order, dephased-control,
and postselected-control arms test equal-prior neutrality, biased-prior
activation, coherence loss under dephasing, and the separation between
joint T+E₂ recovery and E₂-only records. The evidence is organized into
three distinct computational layers.

The record-writing gates used in the circuit validation are

U\_A(α\_A)=\textbar0⟩⟨0\textbar\_S⊗I\_\{E₁\}+\textbar1⟩⟨1\textbar\_S⊗R\_y(α\_A)\_\{E₁\},
(18)

U\_B(θ,α\_B)=R\_y(-θ)\_S{[}\textbar0⟩⟨0\textbar\_S⊗I\_\{E₂\}+\textbar1⟩⟨1\textbar\_S⊗R\_y(α\_B)\_\{E₂\}{]}R\_y(θ)\_S.
(19)

Identity operators on systems not displayed in Eqs. (18) and (19) are
implicit. These definitions match the archived circuit construction used
for the independent gate-level validation.

\includegraphics[width=6.6in,height=2.66037in]{/mnt/data/Causal_Nonseparability_Environmental_Objectivity/figures/pandoc_media/media/image1.png}

\textbf{Figure 1. Four-qubit circuit and access routes for coherent,
dephased, discarded-control, and postselected comparisons. Table VI
summarizes the accessible systems and target prediction for each arm.}

\subsection{A. Circuit-level statevector
validation}\label{a.-circuit-level-statevector-validation}

The archived Qiskit implementation constructs U\_A(α\_A) as a controlled
R\_y record operation from S to E₁ and U\_B(θ,α\_B) by rotating S
through θ, applying the controlled R\_y record operation from S to E₂,
and undoing the rotation. Fixed AB and BA circuits and the coherently
controlled order act on T,S,E₁,E₂. The existing 51-point Qiskit
statevector grid reports a maximum error of 3.33×10\^{}-16 relative to
Eq. (13) and a maximum coherent-minus-best-fixed discrepancy of
4.44×10\^{}-16.

A new NumPy-only gate-level implementation independently constructs the
same 16-dimensional circuit matrices, prepares both binary pointer
labels, evolves all four qubits, explicitly traces to E₂ or T+E₂, and
exports the physical-basis entries of Δ\_AB, Δ\_BA, and Δ\_coh. Across
17 mismatch angles and three strength pairs, the maximum
analytic-formula error and coherent-minus-best-fixed discrepancy are
both 2.22×10\^{}-16. The Version 5 gate-level calculation targets the
equal-prior decision-neutrality theorem. The Holevo benchmark,
many-fragment results, collision-noise surfaces, and local
spectrum-broadcast diagnostics remain supported by their dedicated
archived calculations rather than by this gate-level validation dataset.

\subsection{B. Reduced-normal-form
verification}\label{b.-reduced-normal-form-verification}

The separate reduced-normal-form script begins from the certified
parameter s=sin(α\_B/2)cosθ and constructs valid conditional density
matrices realizing diag(s,-s) and diag(s,-s,0,0). It verifies the
eigenvalues, trace norms, and Helstrom formula on 51 settings. This
calculation checks the consequences of the spectral theorem; it is not
presented as an independent derivation of the four-qubit circuit
spectrum.

\subsection{C. Finite-shot ideal
sampling}\label{c.-finite-shot-ideal-sampling}

A finite-shot calculation samples the ideal projective probabilities at
99 settings with 8,192 shots per pointer label and fixed seed 20260812.
The mean bias is 0.000200, the mean absolute error is 0.002458, and the
maximum absolute error is 0.010897. To avoid the boundary undercoverage
of the earlier plug-in bootstrap, Version 4 uses Bonferroni simultaneous
Clopper-Pearson intervals for the two binomial probabilities and
propagates those intervals through \textbar p₀-p₁\textbar. All 99
analytic targets lie inside the resulting intervals in this validation
grid; this observed grid coverage is a diagnostic, not a universal
coverage claim. The intervals are intentionally conservative, and 100\%
observed inclusion on this finite grid should not be interpreted as
exact nominal calibration.

The finite-shot study isolates ideal sampling fluctuations. It does not
model gate error, decoherence, readout bias, crosstalk, compilation
overhead, or calibration drift.

\subsection{D. Machine-readable
artifacts}\label{d.-machine-readable-artifacts}

The Version 5 validation set includes gate\_level\_validation\_v4.py,
v4\_gate\_level\_validation.csv, v4\_finite\_shot\_validation.csv,
v4\_interval\_noncoverage.csv, v4\_validation\_summary.json,
v4\_checksums.sha256, and the updated validation figure. The gate-level
CSV contains the physical-basis matrix entries, spectra, trace norms,
and Helstrom values for all 51 circuit settings.

\textbf{Table VII. Summary of numerical validation layers.}

\begin{longtable}[]{@{}
  >{\raggedright\arraybackslash}p{(\linewidth - 6\tabcolsep) * \real{0.2500}}
  >{\raggedright\arraybackslash}p{(\linewidth - 6\tabcolsep) * \real{0.2500}}
  >{\raggedright\arraybackslash}p{(\linewidth - 6\tabcolsep) * \real{0.2500}}
  >{\raggedright\arraybackslash}p{(\linewidth - 6\tabcolsep) * \real{0.2500}}@{}}
\toprule\noalign{}
\begin{minipage}[b]{\linewidth}\raggedright
\textbf{Validation layer}
\end{minipage} & \begin{minipage}[b]{\linewidth}\raggedright
\textbf{Settings}
\end{minipage} & \begin{minipage}[b]{\linewidth}\raggedright
\textbf{Primary quantity}
\end{minipage} & \begin{minipage}[b]{\linewidth}\raggedright
\textbf{Maximum discrepancy / error}
\end{minipage} \\
\midrule\noalign{}
\endhead
\bottomrule\noalign{}
\endlastfoot
Qiskit statevector & 51 & Helstrom formula & 3.33×10\^{}-16 \\
NumPy gate level & 51 & Helstrom formula and physical-basis matrices &
2.22×10\^{}-16 \\
Reduced normal form & 51 & Spectral consequence & 0 in stored double
precision \\
Finite-shot ideal sampling & 99 & Helstrom estimate & 0.010897 maximum
absolute error \\
\end{longtable}

\includegraphics[width=6.6in,height=2.49873in]{/mnt/data/Causal_Nonseparability_Environmental_Objectivity/figures/pandoc_media/media/image2.png}

Figure 2. Gate-level and finite-shot validation. Curves show the
analytic prediction; representative error bars are simultaneous
Clopper-Pearson intervals. The right panel displays finite-shot bias
across 99 ideal-sampling settings. Each setting uses 8,192 shots per
pointer label. The vertical guide marks θ=π/2, where the analytic
contrast vanishes. Table VII compares all four numerical-validation
layers.

\textbf{Table VI. Validation arms, accessible systems, and target
predictions.}

\begin{longtable}[]{@{}
  >{\raggedright\arraybackslash}p{(\linewidth - 4\tabcolsep) * \real{0.3333}}
  >{\raggedright\arraybackslash}p{(\linewidth - 4\tabcolsep) * \real{0.3333}}
  >{\raggedright\arraybackslash}p{(\linewidth - 4\tabcolsep) * \real{0.3333}}@{}}
\toprule\noalign{}
\begin{minipage}[b]{\linewidth}\raggedright
\textbf{Arm}
\end{minipage} & \begin{minipage}[b]{\linewidth}\raggedright
\textbf{Accessible systems}
\end{minipage} & \begin{minipage}[b]{\linewidth}\raggedright
\textbf{Primary test}
\end{minipage} \\
\midrule\noalign{}
\endhead
\bottomrule\noalign{}
\endlastfoot
Fixed AB / BA & E₂ or declared fragment & Matched fixed-order
baseline \\
Coherent control & T+E₂ & Joint-access Holevo contribution and Helstrom
neutrality \\
Dephased control & T+E₂ after D\_T & Removal of coherence-sensitive
term \\
Control discarded & E₂ only & Environment-only comparison \\
Postselected control & Outcome Y and E₂ & Probability-weighted
information budget \\
\end{longtable}

\section{X. Scope and limitations}\label{x.-scope-and-limitations}

\begin{itemize}
\item
  The matched-family Holevo advantage is not claimed against the full
  causally separable process cone.
\item
  Causal nonseparability does not imply environmental objectivity, and
  objectivity is not a causal witness.
\item
  The inherited-witness theorem requires deterministic local recovery of
  the certified base process.
\item
  The SBS diagnostics are finite local and correlation diagnostics, not
  a global distance-to-SBS theorem.
\item
  The validation package contains analytic proofs, Qiskit and NumPy
  gate-level statevector results, reduced-normal-form checks, and
  finite-shot ideal-sampling outputs. These are computational validation
  results, not physical-device measurements.
\item
  The access decomposition is defined relative to the declared
  branch-basis dephasing and access structure.
\item
  The Helstrom-neutrality theorem is exact for the declared equal-prior
  binary model and is not asserted for arbitrary tasks, priors,
  channels, or dimensions.
\end{itemize}

\section{XI. Conclusion}\label{xi.-conclusion}

Environmental objectivity and causal nonseparability are independent.
When record-generating operations are coherently ordered, branchwise
pointer compatibility determines deterministic preservation, while the
control provides a bounded information resource. The resource can
increase joint Holevo information without improving equal-prior binary
discrimination or fragment-only redundancy. Postselection is probability
constrained, record replication does not make the bonus extensive, and
deterministic local record outputs preserve the causal certificate
through witness pullback.

The comparison with communication-power results further shows why
coherent order should be assessed task by task: an advantage in one-shot
message transmission does not imply an advantage in equal-prior
environmental-record discrimination. The framework therefore gives
theory and experiment a common checklist for separating redundant
environmental records, control-assisted information, postselection
effects, decision performance, and independent causal certification in
controlled-order record circuits.

\section{Data and code availability}\label{data-and-code-availability}

All deterministic scripts, prediction tables, exact matrix outputs,
finite-shot computational counts, bootstrap intervals, local simulator
outputs, theorem certificates, reports, and checksums assembled for this
submission are included in the accompanying reproducibility package. The
package includes machine-readable numerical validation outputs and
simulation-generated finite-shot data but no measurements from a
physical quantum processor. The final public archive will be cited as:
S. Muniz, Causal Nonseparability and Environmental Objectivity
Reproducibility Package, Version 5.0, Zenodo (2026), DOI:
10.5281/zenodo.XXXXXXX {[}PLACEHOLDER - replace with minted DOI before
submission{]}. See Ref. {[}26{]}.

\section{Author contributions, declarations, and AI
assistance}\label{author-contributions-declarations-and-ai-assistance}

S.M. is the sole author, conceived and directed the research, and takes
responsibility for the manuscript and its claims. The author is
affiliated with Dr. Ware Technology Services and declares no competing
interests. A large language model assisted with drafting, editing,
bibliographic organization, numerical code, symbolic checks, and
critical review. The author is responsible for verification of all
mathematical statements, calculations, references, and final text. No AI
system is an author.

\section{Appendix A. Matrix route for the neutrality
theorem}\label{appendix-a.-matrix-route-for-the-neutrality-theorem}

The original symbolic script generates the physical-basis operators from
the real-rotation circuit and verifies p\_coh(λ)=λ²p\_AB(λ)=λ²p\_BA(λ).
Version 4 independently reconstructs the four-qubit gates with NumPy,
traces to E₂ or T+E₂, and exports the physical-basis entries of Δ\_AB,
Δ\_BA, and Δ\_coh at all 51 grid points in
v4\_gate\_level\_validation.csv. The separate reduced-normal-form file
remains a consequence-level check rather than a circuit-level
derivation.

\section{Appendix B. Referee reproducibility
checklist}\label{appendix-b.-referee-reproducibility-checklist}

\textbf{Table VIII. Claim-to-artifact audit map.}

Table VIII maps each principal claim to its required reproducibility
artifact and verification target.

\begin{longtable}[]{@{}
  >{\raggedright\arraybackslash}p{(\linewidth - 4\tabcolsep) * \real{0.3333}}
  >{\raggedright\arraybackslash}p{(\linewidth - 4\tabcolsep) * \real{0.3333}}
  >{\raggedright\arraybackslash}p{(\linewidth - 4\tabcolsep) * \real{0.3333}}@{}}
\toprule\noalign{}
\begin{minipage}[b]{\linewidth}\raggedright
\textbf{Claim}
\end{minipage} & \begin{minipage}[b]{\linewidth}\raggedright
\textbf{Required artifact}
\end{minipage} & \begin{minipage}[b]{\linewidth}\raggedright
\textbf{Verification target}
\end{minipage} \\
\midrule\noalign{}
\endhead
\bottomrule\noalign{}
\endlastfoot
Access decomposition & Symbolic notebook or script & Eq. (3) and
marginal-invariance check \\
Postselection budget & Analytic derivation and sampled instruments &
Eqs. (7)-(9) \\
Helstrom neutrality & Symbolic determinant certificate and statevector
grid & Eqs. (10)-(13) \\
Matched Holevo bonus & Parameter file and machine-readable output &
0.0839771759-bit difference \\
Many-fragment result & Subset-enumeration script and tables & N=2,4,6,8
conclusions \\
Witness pullback & Process-space assumptions and witness certificate &
Eqs. (14)-(17) \\
Four-qubit route & OpenQASM and tomography scripts & Declared simulator
predictions \\
Gate-level circuit neutrality & gate\_level\_validation\_v4.py;
v4\_gate\_level\_validation.csv & Physical-basis matrices, spectra, and
Eq. (13) over 51 settings \\
Reduced spectral normal forms & classical\_validation\_v3.py;
v3\_exact\_matrix\_validation.csv & Normal-form eigenvalues and trace
norms over 51 settings \\
Finite-shot ideal sampling & v4\_finite\_shot\_validation.csv & Counts,
bias, error, and simultaneous intervals over 99 settings \\
Validation summary and integrity & v4\_validation\_summary.json;
v4\_checksums.sha256 & Seed, versions, aggregate metrics, and SHA-256
hashes \\
Version 5 manuscript validation & Submission\_Paper\_PRXQ\_V5.docx;
v4\_checksums.sha256 & Gate definitions, tensor ordering,
cross-references, and validated numerical summaries \\
\end{longtable}

\section{References}\label{references}

{[}1{]} W. H. Zurek, Decoherence, einselection, and the quantum origins
of the classical, Rev. Mod. Phys. 75, 715 (2003).

{[}2{]} H. Ollivier, D. Poulin, and W. H. Zurek, Objective Properties
from Subjective Quantum States: Environment as a Witness, Phys. Rev.
Lett. 93, 220401 (2004).

{[}3{]} R. Blume-Kohout and W. H. Zurek, Quantum Darwinism:
Entanglement, branches, and the emergent classicality of redundantly
stored quantum information, Phys. Rev. A 73, 062310 (2006).

{[}4{]} W. H. Zurek, Quantum Darwinism, Nat. Phys. 5, 181 (2009).

{[}5{]} C. J. Riedel and W. H. Zurek, Quantum Darwinism in an Everyday
Environment, Phys. Rev. Lett. 105, 020404 (2010).

{[}6{]} J. K. Korbicz, P. Horodecki, and R. Horodecki, Objectivity in a
Noisy Photonic Environment through Quantum State Information
Broadcasting, Phys. Rev. Lett. 112, 120402 (2014).

{[}7{]} R. Horodecki, J. K. Korbicz, and P. Horodecki, Quantum origins
of objectivity, Phys. Rev. A 91, 032122 (2015).

{[}8{]} J. K. Korbicz, Roads to objectivity: Quantum Darwinism, spectrum
broadcast structures, and strong quantum Darwinism, Quantum 5, 571
(2021).

{[}9{]} O. Oreshkov, F. Costa, and Č. Brukner, Quantum correlations with
no causal order, Nat. Commun. 3, 1092 (2012).

{[}10{]} M. Araújo et al., Witnessing causal nonseparability, New J.
Phys. 17, 102001 (2015).

{[}11{]} C. Branciard, Witnesses of causal nonseparability: an
introduction and a few case studies, Sci. Rep. 6, 26018 (2016).

{[}12{]} J. Wechs, A. A. Abbott, and C. Branciard, Multipartite causal
(non)separability, New J. Phys. 21, 013027 (2019).

{[}13{]} E. Castro-Ruiz, F. Giacomini, and Č. Brukner, Dynamics of
quantum causal structures, Phys. Rev. X 8, 011047 (2018).

{[}14{]} S. Campbell et al., Collisional unfolding of quantum Darwinism,
Phys. Rev. A 99, 042103 (2019).

{[}15{]} M. Antesberger et al., Higher-Order Process Matrix Tomography
of a Passively-Stable Quantum Switch, PRX Quantum 5, 010325 (2024).

{[}16{]} M. M. Taddei, R. V. Nery, and L. Aolita, Quantum superpositions
of causal orders as an operational resource, Phys. Rev. Research 1,
033174 (2019).

{[}17{]} E. Doucet and S. Deffner, Classifying Two-Body Hamiltonians for
Quantum Darwinism, Phys. Rev. X 14, 041064 (2024).

{[}18{]} D. A. Chisholm, G. M. Palma, and L. Innocenti, Emergence of
quantum Darwinism and pointer states for noncommuting evolutions, APS
Open Science 1, 000024 (2026).

{[}19{]} G. Chiribella, G. M. D\textquotesingle Ariano, P. Perinotti,
and B. Valiron, Quantum computations without definite causal structure,
Phys. Rev. A 88, 022318 (2013). DOI: 10.1103/PhysRevA.88.022318.

{[}20{]} T. P. Le and A. Olaya-Castro, Strong Quantum Darwinism and
Strong Independence are Equivalent to Spectrum Broadcast Structure,
Phys. Rev. Lett. 122, 010403 (2019). DOI:
10.1103/PhysRevLett.122.010403.

{[}21{]} L. A. Rozema, T. Strömberg, H. Cao, Y. Guo, B.-H. Liu, and P.
Walther, Experimental aspects of indefinite causal order in quantum
mechanics, Nat. Rev. Phys. 6, 483-499 (2024). DOI:
10.1038/s42254-024-00739-8.

{[}22{]} X. Zhao, B. Zhao, C. Branciard, and G. Chiribella, The
communication power of indefinite causal order, arXiv:2510.08507v2
(2026). DOI: 10.48550/arXiv.2510.08507.

{[}23{]} C. M. D. Richter, M. Antesberger, H. Cao, P. Walther, and L. A.
Rozema, Toward an Experimental Device-Independent Verification of
Indefinite Causal Order, PRX Quantum 7, 010354 (2026). DOI:
10.1103/5t2y-ddmt.

{[}24{]} Z. Zhu et al., Observation of quantum Darwinism and the origin
of classicality with superconducting circuits, Sci. Adv. 11, eadx6857
(2025). DOI: 10.1126/sciadv.adx6857.

{[}25{]} A. Steffinlongo and H. Dourdent, Simulating Noncausality with
Quantum Control of Causal Orders, Phys. Rev. Research 8, 013127 (2026).
DOI: 10.1103/9jss-yf1c.

{[}26{]} S. Muniz, Causal Nonseparability and Environmental Objectivity
Reproducibility Package, Version 5.0, Zenodo (2026), DOI:
10.5281/zenodo.XXXXXXX {[}PLACEHOLDER{]}.

\end{document}
